\begin{table}[!htbp]
\centering
\caption{Paired effect sizes and empirical quartiles at the declared statistical unit.}
\label{tab:effect_sizes}
\small
\setlength{\tabcolsep}{4pt}
\begin{tabular}{l S[table-format=2.0] *{3}{S[table-format=2.3]} *{3}{S[table-format=5.1]}}
\toprule
{Cohort} & {$n$} & {$D_{25}$} & {$D_{50}$} & {$D_{75}$} & {$R_{25}$} & {$R_{50}$} & {$R_{75}$} \\
\midrule
Burgers & 12 & 24.814 & 27.636 & 29.749 & 282.8 & 335.4 & 574.1 \\
Allen--Cahn & 9 & 19.158 & 20.006 & 21.633 & 56.9 & 78.6 & 81.2 \\
Saturation & 6 & 16.532 & 16.815 & 17.753 & 10427.2 & 16883.1 & 18262.9 \\
\bottomrule
\end{tabular}
\par\vspace{3pt}\begin{minipage}{.98\linewidth}\footnotesize $D=E_{\rm raw}-E_{\rm SAEPS}$ and $R=E_{\rm raw}/E_{\rm SAEPS}$ use the same unit-level pair. Subscripts denote empirical percentiles, using linear interpolation. Historical units are accepted checkpoints; saturation units are six four-fit medians. The saturation difference is formed after that pooling. These quartiles summarize dispersion, not a bootstrap or calibrated 95\% confidence interval. No cohorts are pooled for inference.\end{minipage}
\end{table}
